\paperchapter{December 2013}{December 2013}{A}{60}{GSET 2013 (II)}
\keytable{\kcell{1}{A} & \kcell{2}{D} & \kcell{3}{C} & \kcell{4}{C} & \kcell{5}{C} & \kcell{6}{B} & \kcell{7}{D} & \kcell{8}{A} & \kcell{9}{C} & \kcell{10}{B} \\
\kcell{11}{D} & \kcell{12}{C} & \kcell{13}{D} & \kcell{14}{C} & \kcell{15}{B} & \kcell{16}{D} & \kcell{17}{D} & \kcell{18}{D} & \kcell{19}{D} & \kcell{20}{C} \\
\kcell{21}{C} & \kcell{22}{C} & \kcell{23}{C} & \kcell{24}{D} & \kcell{25}{B} & \kcell{26}{A} & \kcell{27}{C} & \kcell{28}{D} & \kcell{29}{B} & \kcell{30}{C} \\
\kcell{31}{D} & \kcell{32}{A} & \kcell{33}{C} & \kcell{34}{B} & \kcell{35}{B} & \kcell{36}{A} & \kcell{37}{D} & \kcell{38}{C} & \kcell{39}{C} & \kcell{40}{B} \\
\kcell{41}{B} & \kcell{42}{A} & \kcell{43}{B} & \kcell{44}{C} & \kcell{45}{D} & \kcell{46}{C} & \kcell{47}{D} & \kcell{48}{B} & \kcell{49}{B} & \kcell{50}{C} \\
\kcell{51}{D} & \kcell{52}{B} & \kcell{53}{A} & \kcell{54}{A} & \kcell{55}{C} & \kcell{56}{B} & \kcell{57}{C} & \kcell{58}{C} & \kcell{59}{B} & \kcell{60}{A}}
\begin{question}{1}
{Learning takes place when:\par}
\optsone{there is a match between `what' and `how' students learn}{the teacher presents the content accurately}{the learning experience is organised within a stipulated time}{teacher's focus is on completion of syllabus}
\answer{A}{there is a match between `what' and `how' students learn}
\begin{explain}
Learning is most effective when the \textbf{content (``what'') matches the method and style (``how'')} by which students learn, i.e.\ when teaching suits the learner's needs, level and learning style.
\why{Accurate presentation, time limits or completing the syllabus are teacher-centred concerns and do not by themselves guarantee learning.}
\end{explain}
\end{question}

\begin{question}{2}
{Which of the following does not justify teaching as a profession?\par}
\optstwo{Specialized knowledge}{Specialized skills}{Marked Training Institution}{Market Oriented effort}
\answer{D}{Market Oriented effort}
\begin{explain}
A profession is marked by specialised knowledge and skills, formal training institutions, a code of ethics and a service motive. A \textbf{market-oriented (profit-seeking) effort} is a feature of business, not a justification of teaching as a profession.
\end{explain}
\end{question}

\begin{question}{3}
{Which is not a significant challenge faced by a beginner teacher?\par}
\optstwo{Content Mastery}{Use of ICT in teaching}{Following policy of the institution}{Classroom management}
\answer{C}{Following policy of the institution}
\begin{explain}
Beginning teachers typically struggle with content mastery, use of ICT and classroom management (Veenman's study of novice teachers). \textbf{Following institutional policy} is routine and is not a significant challenge.
\end{explain}
\end{question}

\begin{question}{4}
{It has been usually opined that `learning in the educational institutions these days is wasted' because:\par}
\optsone{subjects are boring}{teachers communicate their subject thoroughly}{subjects are not communicated in life relevant ways}{most teachers are well acquainted with their subject}
\answer{C}{subjects are not communicated in life relevant ways}
\begin{explain}
Learning is said to be ``wasted'' when subjects are taught in an abstract, bookish way that is \textbf{not related to real life}. Students then cannot see the relevance of what they learn and do not retain or use it.
\end{explain}
\end{question}

\begin{question}{5}
{A teacher uses local language to raise the effectiveness in class. This refers to:\par}
\optstwo{effective encoding}{effective decoding}{effective encoding and decoding}{effective media interlink}
\answer{C}{effective encoding and decoding}
\begin{explain}
Using the local language makes the message understandable both when the teacher \textbf{encodes} it (puts ideas in familiar words) and when students \textbf{decode} it (understand it easily). It improves \textbf{encoding and decoding} together.
\end{explain}
\end{question}

\begin{question}{6}
{Which of the following ensures greater participation of students in classroom?\par}
\optstwo{Lecture by the teacher}{Discussion organised by the teacher}{Field trip organised by the teacher}{Class notes given by the teacher}
\answer{B}{Discussion organised by the teacher}
\begin{explain}
A teacher-organised \textbf{discussion} involves all students in speaking, listening, questioning and reasoning, so it ensures the greatest participation.
\why{Lecture and class notes are one-way; a field trip is occasional and not a routine classroom activity.}
\end{explain}
\end{question}

\begin{question}{7}
{A characteristic to be ascertained initially in a good research is:\par}
\optstwo{Well formulated hypothesis}{A good research supervisor}{Adequate library work}{A well formulated problem}
\answer{D}{A well formulated problem}
\begin{explain}
Research begins with a \textbf{well-formulated problem}; hypotheses, literature review and design all follow from it. ``A problem well stated is half solved'' (Dewey).
\end{explain}
\end{question}

\begin{question}{8}
{You want to study the opinion of students of Gujarat regarding introduction of `choice based credit system' in higher education. The most appropriate sampling technique for the purpose is:\par}
\optstwo{Multistage Random Sampling}{Cluster Sampling}{Quota Sampling}{Snowball Sampling}
\answer{A}{Multistage Random Sampling}
\begin{explain}
Students across Gujarat form a large, geographically spread population. \textbf{Multistage random sampling} (e.g.\ districts $\rightarrow$ universities $\rightarrow$ colleges $\rightarrow$ students) gives a representative sample efficiently.
\end{explain}
\end{question}

\begin{question}{9}
{The research that generates new ideas, principles and theories providing foundation for progress and development in different fields is called:\par}
\optstwo{Empirical Research}{Historical Research}{Fundamental Research}{Experimental Research}
\answer{C}{Fundamental Research}
\begin{explain}
\textbf{Fundamental (basic) research} generates new knowledge, principles and theories, forming the foundation for later applied work.
\end{explain}
\end{question}

\begin{question}{10}
{Which of the following is not the basis of scientific method?\par}
\optstwo{Empirical evidence}{Opinions and responses of people}{Laboratory experimentations}{Objective considerations}
\answer{B}{Opinions and responses of people}
\begin{explain}
The scientific method rests on \textbf{empirical evidence}, \textbf{experimentation} and \textbf{objectivity}. Mere \textbf{opinions and responses of people} are subjective and are not the basis of the scientific method.
\end{explain}
\end{question}

\begin{question}{11}
{Researcher carries out Review of Related Literature. This helps her in many ways. But one of the following is not covered by the review of related literature:\par}
\optstwo{finding gaps in research}{getting research questions}{framing hypothesis}{implications of research}
\answer{D}{implications of research}
\begin{explain}
The literature review helps identify research gaps, frame research questions and formulate hypotheses. The \textbf{implications of research} come from the researcher's \emph{own findings} at the end of the study, not from the review.
\end{explain}
\end{question}

\begin{question}{12}
{\blank{} is unsuitable for graphical representation of continuous data.\par}
\optsfour{Histogram}{Polygon}{Bar Chart}{Ogive}
\answer{C}{Bar Chart}
\begin{explain}
Continuous data are shown with a histogram, frequency polygon or ogive (cumulative frequency curve). A \textbf{bar chart}, with gaps between bars, suits \emph{discrete or categorical} data and is unsuitable for continuous data.
\end{explain}
\end{question}

\begin{question}{13}
{Which of the following arrangements from the smallest to the largest unit of storage holds good in information technology?\par}
\optstwo{Kilobyte, Gigabyte, Megabyte, Terabyte}{Kilobyte, Terabyte, Megabyte, Gigabyte}{Kilobyte, Terabyte, Gigabyte, Megabyte}{Kilobyte, Megabyte, Gigabyte, Terabyte}
\answer{D}{Kilobyte, Megabyte, Gigabyte, Terabyte}
\begin{explain}
From smallest to largest: \textbf{KB $<$ MB $<$ GB $<$ TB}.
\end{explain}
\end{question}

\begin{question}{14}
{Which one of the following is not a network device?\par}
\optsfour{Router}{Switch}{CPU}{Hub}
\answer{C}{CPU}
\begin{explain}
Routers, switches and hubs are \textbf{networking devices}. The \textbf{CPU} is the processing unit of a computer, not a network device.
\end{explain}
\end{question}

\begin{question}{15}
{Universities in India can launch their own FM radio stations. The objective is to:\par}
\optstwo{entertain students}{expose students to further information}{provide students chances to compare}{provide students chances to sing}
\answer{B}{expose students to further information}
\begin{explain}
Campus community radio (licensed since 2004) aims at educational and developmental broadcasting, i.e.\ \textbf{exposing students and the community to further information}.
\end{explain}
\end{question}

\begin{question}{16}
{Find the odd one out:\par}
\optsfour{Windows}{Linux}{Mac}{Google}
\answer{D}{Google}
\begin{explain}
Windows, Linux and macOS are \textbf{operating systems}. \textbf{Google} is a company and search engine, so it is the odd one out. (Google's operating systems are Android and ChromeOS.)
\end{explain}
\end{question}

\begin{question}{17}
{Genuine communication cannot be planned because:\par}
\optsone{many channels need to be utilized}{it is difficult to stick to the pre-decided topic}{language is a very poor medium of expression}{the cycle of messages depends on spontaneous, verbal and non-verbal feedback chain}
\answer{D}{the cycle of messages depends on spontaneous, verbal and non-verbal feedback chain}
\begin{explain}
Genuine (two-way) communication is dynamic. Each message depends on the \textbf{spontaneous verbal and non-verbal feedback} of the other person, so it cannot be fully planned in advance.
\end{explain}
\end{question}

\begin{question}{18}
{The full form of `http' is:\par}
\optstwo{hypertext technology program}{hypertext transfer program}{hypertext topology protocol}{hypertext transfer protocol}
\answer{D}{hypertext transfer protocol}
\begin{explain}
\textbf{HTTP} is \textbf{HyperText Transfer Protocol}, the protocol of the World Wide Web.
\end{explain}
\end{question}

\begin{question}{19}
{Communication is a process of transferring information from one source to other individual, in which:\par}
\optsone{sender should understand the information under transfer}{receiver should understand the information under transfer}{sender sends the information to receiver}{both should understand the information under transfer}
\answer{D}{both should understand the information under transfer}
\begin{explain}
Communication is successful only when there is shared meaning: \textbf{both sender and receiver understand} the information being transferred.
\end{explain}
\end{question}

\begin{question}{20}
{Video Conferencing can be classified as one of the following types of communication:\par}
\optstwo{Visual one way}{Audio-Visual one way}{Audio-Visual two way}{Visual two way}
\answer{C}{Audio-Visual two way}
\begin{explain}
Video conferencing carries \textbf{sound and pictures} in \textbf{both directions} in real time, so it is audio-visual two-way communication.
\end{explain}
\end{question}

\begin{question}{21}
{Voice over Internet Protocol (VoIP) is the use of internet protocol networks to carry voice phone calls. Which of the following would not normally be regarded as an advantage of VoIP?\par}
\optstwo{Low cost telephony}{Increased portability}{Increased geographical anonymity}{Greater functionality}
\answer{C}{Increased geographical anonymity}
\begin{explain}
VoIP offers low-cost calls, portability (one number anywhere) and extra features. \textbf{``Increased geographical anonymity''} is not normally regarded as an advantage; it complicates emergency services and legal tracing.
\end{explain}
\end{question}

\begin{question}{22}
{The process of understanding information within oneself is called:\par}
\optstwo{Interpersonal communication}{Inter-group communication}{Intrapersonal communication}{Interactive communication}
\answer{C}{Intrapersonal communication}
\begin{explain}
Communication within oneself (thinking, self-talk, reflection) is \textbf{intrapersonal communication}.
\end{explain}
\end{question}

\begin{question}{23}
{`Bluetooth' technology allows:\par}
\optsone{Landline phone to mobile phone communication}{Signal transmission on mobile phone only}{Wireless communication between equipments}{Landline transmission only}
\answer{C}{Wireless communication between equipments}
\begin{explain}
\textbf{Bluetooth} is a short-range \textbf{wireless} technology (2.4 GHz band, about 10 m) for connecting devices such as phones, headsets, keyboards and speakers.
\end{explain}
\end{question}

\begin{question}{24}
{Communication will be effective if:\par}
\optsone{it is delivered slowly and clearly}{it is delivered in a calm situation}{it reaches the receiver completely}{it reaches the receiver as intended by the sender}
\answer{D}{it reaches the receiver as intended by the sender}
\begin{explain}
Communication is effective when the message \textbf{reaches the receiver as intended by the sender}, i.e.\ the intended meaning is understood.
\end{explain}
\end{question}

\begin{question}{25}
{Insert the missing number in the following sequence:\newline 5, 11, 23, \blank{}, 95, 191.\par}
\optsfour{57}{47}{44}{49}
\answer{B}{47}
\begin{explain}
Each term is (previous $\times2$) $+1$: $5\rightarrow11\rightarrow23\rightarrow\mathbf{47}\rightarrow95\rightarrow191$.
\end{explain}
\end{question}

\begin{question}{26}
{Letters in the set ARUTF : DSXUI have certain relationship. Under the same relationship, mark the correct choice for the letters of the set PQRSF : \blank{}.\par}
\optsfour{SRUTI}{SRUVI}{RSUTI}{RSUVG}
\answer{A}{SRUTI}
\begin{explain}
ARUTF $\rightarrow$ DSXUI: the shifts are $+3, +1, +3, +1, +3$.
PQRSF $\rightarrow$ P$+3=$S, Q$+1=$R, R$+3=$U, S$+1=$T, F$+3=$I, giving \textbf{SRUTI}.
\end{explain}
\end{question}

\begin{question}{27}
{In a certain code RSDFRBMUG means SREESANTH, then how FIXING can be written in the same code?\par}
\optsfour{EJWHPF}{GHYHOF}{EJWJMH}{GHWJMH}
\answer{C}{EJWJMH}
\begin{explain}
SREESANTH $\rightarrow$ RSDFRBMUG: the shifts alternate $-1, +1, -1, +1, \ldots$
FIXING $\rightarrow$ F$-1=$E, I$+1=$J, X$-1=$W, I$+1=$J, N$-1=$M, G$+1=$H, giving \textbf{EJWJMH}.
\end{explain}
\end{question}

\begin{question}{28}
{At what time between 4 and 5 o'clock the minute and hour hands of a clock point in opposite directions?\par}
\optstwo{45 minutes past 4}{40 minutes past 4}{$50\frac{4}{11}$ minutes past 4}{$54\frac{6}{11}$ minutes past 4}
\answer{D}{$54\frac{6}{11}$ minutes past 4}
\begin{explain}
The hands are opposite when they are $180^\circ$ apart. The angle is $|30H-5.5M|$, so $|120-5.5M|=180$ gives $5.5M=300$ and
\[M=\frac{300}{5.5}=\frac{600}{11}=54\tfrac{6}{11}\text{ minutes past 4}\]
\end{explain}
\end{question}

\begin{question}{29}
{If PEN = 25 and HEN = 17, then MEN = ?\par}
\optsfour{20}{22}{32}{30}
\answer{B}{22}
\begin{explain}
PEN and HEN differ only in the first letter: P $=16$, H $=8$, and $25-17=8$ matches $16-8$. So the value is (position of first letter) $+9$: M $=13$, giving $13+9=\mathbf{22}$.
\end{explain}
\end{question}

\begin{question}{30}
{In a row of girls, Bindu is twelfth from either end of the row. How many girls are there in the row?\par}
\optsfour{13}{24}{23}{25}
\answer{C}{23}
\begin{explain}
Bindu is 12th from each end: $11$ girls on each side $+1=\mathbf{23}$.
\end{explain}
\end{question}

\begin{context}
\textbf{Read the following text and answer the questions from 31 to 36.}\par\smallskip
Investigators have noticed some similarities between the fantastic events in fairy tales and those in adult dreams --- the fulfilment of wishes, winning over all competitors and the destruction of enemies. There are, of course, very significant differences also. For example, in dreams more often than not the wish fulfilment is disguised, while in fairy tales much of it is openly expressed. To a considerable degree, dreams are the result of inner pressures which have found no relief, and of problems which beset a person to which he knows no solution and to which the dream finds none. The fairy tale does the opposite; it projects the relief of all pressures and not only offers ways to solve problems, but promises a happy solution. We cannot control what goes on in our dreams. Although our inner censorship influences what we may dream, such control occurs at an unconscious level. The fairy tale, on the other hand, is the result of conscious and unconscious content having been shaped by the conscious mind, not of one particular person but of many in regard to what they view as universal human problems. Only if a fairy tale met the conscious and unconscious needs of many people, would it be told and retold with great interest.
\end{context}
\begin{question}{31}
{It can be inferred from the passage that the author's interest in fairy tales centres chiefly on their:\par}
\optstwo{literary qualities}{historical background}{factual accuracy}{psychological relevance}
\answer{D}{psychological relevance}
\begin{explain}
The author discusses wishes, inner pressures, the conscious and the unconscious, i.e.\ the \textbf{psychological relevance} of fairy tales (the passage echoes Bruno Bettelheim's \emph{The Uses of Enchantment}).
\end{explain}
\end{question}

\begin{question}{32}
{Which of the following is common in dreams and fairy tales?\par}
\optstwo{Destruction of enemies}{Solution of universal human problems}{Unconscious requirement of many people}{Subject to inner censorship}
\answer{A}{Destruction of enemies}
\begin{explain}
The similarities listed are ``the fulfilment of wishes, winning over all competitors and the \textbf{destruction of enemies}''.
\end{explain}
\end{question}

\begin{question}{33}
{According to the passage, fairy tales differ from dreams in which of the following characteristics?\newline (I) The communal nature of their creation\newline (II) Their convention of a happy ending\newline (III) Their enduring general appeal.\par}
\optstwo{(I) only}{(II) only}{(I) and (II) only}{(II) and (III) only}
\answer{C}{(I) and (II) only}
\begin{explain}
Fairy tales differ from dreams in being shaped by \textbf{many people} (communal creation, I) and in promising a \textbf{happy solution} (II). Enduring appeal (III) is a consequence, not a stated point of difference. The answer is (I) and (II) only.
\end{explain}
\end{question}

\begin{question}{34}
{Which of the following is not true?\par}
\optsone{Fairy tales offer solutions to human problems}{All tales are retold with great interest}{Dreams and fairy tales differ in some ways}{Dreams are subject to inner censorship}
\answer{B}{All tales are retold with great interest}
\begin{explain}
The passage says only a tale that meets the needs of many people is retold with great interest, i.e.\ \emph{not all} tales. So ``all tales are retold with great interest'' is \textbf{not true}.
\end{explain}
\end{question}

\begin{question}{35}
{In the context of this passage, the relationship between dreams and fairy tales is:\par}
\optsone{cause and effect}{comparison and contrast}{general principle and example}{similar examples of a general principle}
\answer{B}{comparison and contrast}
\begin{explain}
The passage sets out similarities and then differences between dreams and fairy tales: \textbf{comparison and contrast}.
\end{explain}
\end{question}

\begin{question}{36}
{An appropriate title for the passage would be:\par}
\optsone{Psychology of dreams}{Dreams and fairy tales}{How the Unconscious Manifests}{Solve your problems through dreams and stories}
\answer{A}{Psychology of dreams}
\keynote{Many readers would find ``Dreams and fairy tales'' (B) a closer title, since the passage compares both throughout. The official key is (A).}
\begin{explain}
The official key gives \textbf{Psychology of dreams}: the passage analyses wish fulfilment, inner pressures, censorship and the conscious and unconscious mind, which are psychological aspects of dreams (with fairy tales used for contrast).
\end{explain}
\end{question}

\begin{question}{37}
{Cigarette smoking is banned in public places in India because it:\par}
\optstwo{Irritates non-smokers}{Suffocates non-smokers}{Creates headache to non-smokers}{Harms lungs of non-smokers}
\answer{D}{Harms lungs of non-smokers}
\begin{explain}
Passive (second-hand) smoke \textbf{harms the lungs and health of non-smokers}. This health hazard is the basis of the ban under the Cigarettes and Other Tobacco Products Act (COTPA), 2003, and the 2008 rules on smoking in public places.
\end{explain}
\end{question}

\begin{question}{38}
{Waste management has been one of the most current issues being tackled at present. Which of the following is a good example for it?\par}
\optstwo{Keep parks and roads clean}{Use less water to clean vessels}{Recycle biodegradable waste}{Reduce use of electronic gadgets}
\answer{C}{Recycle biodegradable waste}
\begin{explain}
Waste management follows ``reduce, reuse, \textbf{recycle}''. \textbf{Recycling biodegradable waste} (composting, biogas) is a direct example.
\end{explain}
\end{question}

\begin{question}{39}
{A few actions of youth that point to their environmental insensitivity are given below. Which does not belong to this category?\par}
\optstwo{grafitti on the walls and desks}{playing audio loud}{going on strike}{riding motorbikes without silencers}
\answer{C}{going on strike}
\begin{explain}
Graffiti, loud audio and noisy motorbikes harm the environment (visual and noise pollution). \textbf{Going on strike} is a social or political action, not an environmental one.
\end{explain}
\end{question}

\begin{question}{40}
{Promotion of non-conventional sources of energy is advocated in day to day life now-a-days. This will reduce:\par}
\optsone{cost of petrol and diesel}{consumption of crude oil}{dependency on electricity and petroleum based fuels}{cost of electricity}
\answer{B}{consumption of crude oil}
\begin{explain}
Non-conventional sources (solar, wind, biogas) replace fossil fuels and so reduce the \textbf{consumption of crude oil}.
\end{explain}
\end{question}

\begin{question}{41}
{"Narmada Bachao Andolan" was triggered due to:\par}
\optsone{disproportionate sharing of Narmada water with neighbouring states}{disturbances and rehabilitation issues}{water pollution to the farming land}{ill-treatment to farmers}
\answer{B}{disturbances and rehabilitation issues}
\begin{explain}
The \textbf{Narmada Bachao Andolan} (Medha Patkar, Baba Amte) arose from the \textbf{displacement and rehabilitation} of people affected by the Sardar Sarovar and other Narmada dams, along with environmental concerns.
\end{explain}
\end{question}

\begin{question}{42}
{Which of the following is the cause of change in seasons?\par}
\optsone{The tilt of the earth's axis}{The rotation of the earth}{Distance of the place from equator}{Prevailing winds blowing across land and water}
\answer{A}{The tilt of the earth's axis}
\begin{explain}
Seasons are caused by the \textbf{tilt of the Earth's axis} (about $23.5^\circ$) as the Earth revolves around the Sun. This changes the angle of sunlight and the length of the day.
\why{Rotation causes day and night.}
\end{explain}
\end{question}

\begin{question}{43}
{Evaluate the following statements in terms of whether each is a Fact, Opinion, Prejudice or Advice by choosing correct alternative:\newline (i) Union elections in universities should not be held.\newline (ii) Student leaders are never able to perform well in the study.\newline (iii) Do study for few hours a day regularly even if exams are far away.\newline (iv) Unless necessary credits are earned, degree cannot be awarded.\newline Options are in the order: Opinion, Prejudice, Fact, Advice.\par}
\optstwo{(iii), (i), (ii), (iv)}{(i), (ii), (iv), (iii)}{(i), (ii), (iii), (iv)}{(iv), (i), (ii), (iii)}
\answer{B}{(i), (ii), (iv), (iii)}
\begin{explain}
(i) ``Union elections should not be held'' is an \textbf{opinion}. (ii) ``Student leaders are never able to perform well'' is a \textbf{prejudice} (an unfair generalisation). (iii) ``Do study regularly'' is \textbf{advice}. (iv) ``Degree cannot be awarded without credits'' is a \textbf{fact}.
Opinion, Prejudice, Fact, Advice $=$ (i), (ii), (iv), (iii), which is option (B).
\end{explain}
\end{question}

\begin{question}{44}
{Which of the following diagrams indicates the best relation between Travellers, Train and Bus?\par}
\optsfour{\tikz[baseline=-0.5ex,scale=0.3]{\draw (0,0) circle(1.2); \draw (0,0) circle(0.6); \draw (2.2,0) circle(0.7);}}{\tikz[baseline=-0.5ex,scale=0.3]{\draw (0,0) circle(1); \draw (2.2,0) circle(0.75); \draw (3.8,0) circle(0.5);}}{\tikz[baseline=-0.5ex,scale=0.3]{\draw (0,0) circle(1.1); \draw (-1.2,0.15) circle(0.6); \draw (1.2,0.15) circle(0.6);}}{\tikz[baseline=-0.5ex,scale=0.3]{\draw (0,0) circle(1.3); \draw (-0.55,0) circle(0.5); \draw (0.55,0) circle(0.5);}}
\answer{C}{\tikz[baseline=-0.5ex,scale=0.3]{\draw (0,0) circle(1.1); \draw (-1.2,0.15) circle(0.6); \draw (1.2,0.15) circle(0.6);}}
\begin{explain}
The official key gives diagram (C): a large circle (\textbf{travellers}) partly overlapping two separate circles (\textbf{train} and \textbf{bus}). Some travellers use the train and some the bus, while the train and bus sets also contain elements that are not travellers (crew, vehicles). No one is both in a train and in a bus at the same time.
\tip{Compare GSET Nov 2022 Q.~45, where the key showed train and bus inside the travellers circle. Read each figure carefully.}
\end{explain}
\end{question}

\begin{question}{45}
{Of the following statements there are two statements both of which cannot be true but both can be false. Which are these two statements?\newline (i) All the cricketers are match fixers.\newline (ii) Some cricketers are involved in match fixing\newline (iii) No cricketer is involved in match fixing\newline (iv) Some cricketers are dedicated to cricket.\par}
\optsfour{(iii) and (iv)}{(i) and (ii)}{(ii) and (iii)}{(i) and (iv)}
\answer{D}{(i) and (iv)}
\keynote{Strictly, the pair that cannot both be true but can both be false (\emph{contraries}) is (i) and (iii), which is not offered. The official key gives (D).}
\begin{explain}
Contrary propositions (``All S are P'' and ``No S is P'') cannot both be true but can both be false. Here those are (i) and (iii). (ii) and (iii) are \emph{contradictories} (exactly one is true). The key's answer (D) pairs (i) with (iv), and these two statements are about different predicates. Mark (D) as per the key.
\end{explain}
\end{question}

\begin{question}{46}
{Consider the following statement followed by two assumptions (i) and (ii). Choose the correct answer.\newline Statement: Success is how much a person bounces up after hitting the bottom.\newline Assumptions:\newline (i) A person can totally fail.\newline (ii) It is possible for a person to rise again after failure.\par}
\optsone{Only assumption (i) is implicit.}{Only assumption (ii) is implicit.}{Both the assumptions (i) and (ii) are implicit.}{None of the assumptions (i) and (ii) is implicit.}
\answer{C}{Both the assumptions (i) and (ii) are implicit.}
\begin{explain}
``Bouncing up after hitting the bottom'' assumes that a person can hit the bottom, i.e.\ fail completely (i), and can rise again after failure (ii). \textbf{Both} assumptions are implicit.
\end{explain}
\end{question}

\begin{question}{47}
{In the following, a statement is followed by two conclusions (i) and (ii). Decide which conclusion is valid.\newline Statement: Out of 200 runs made by a team in a one day cricket match, 160 runs were made by spinners.\newline Conclusions:\newline (i) 80\% of the team consists of spinners.\newline (ii) The opening batsmen were spinners.\par}
\optsone{Only conclusion (i) follows.}{Only conclusion (ii) follows.}{Either conclusion (i) or conclusion (ii) follows.}{Neither conclusion (i) nor conclusion (ii) follows.}
\answer{D}{Neither conclusion (i) nor conclusion (ii) follows.}
\begin{explain}
Runs scored say nothing about how many players were spinners (i) or who opened the batting (ii). \textbf{Neither} conclusion follows.
\end{explain}
\end{question}

\begin{question}{48}
{P, Q, R and S are statements such that if P is true, then Q is true and if P is false, then Q is false. The statements P and Q are:\par}
\optsfour{Independent}{Equivalent}{Sub-contrary}{Contradictory}
\answer{B}{Equivalent}
\begin{explain}
P and Q always have the same truth value (P true $\Leftrightarrow$ Q true), so they are \textbf{logically equivalent}.
\end{explain}
\end{question}

\begin{question}{49}
{A magazine printed a survey in its monthly issue and asked readers to fill it out and send it in. Over 1000 students did so. This type of sample is called:\par}
\optstwo{A cluster sample}{A self-selected sample}{A stratified sample}{A simple random sample}
\answer{B}{A self-selected sample}
\begin{explain}
Readers who choose to respond select themselves, so this is a \textbf{self-selected (voluntary response)} sample, which is prone to bias.
\end{explain}
\end{question}

\begin{context}
A survey asked people how often they exceed speed limits. The data is then categorized into the following table of counts showing the relationship between age group and response. Read the table and answer the Question Nos.~50 to 54.
\begin{center}
\begin{tabular}{lccc}\toprule
& \multicolumn{2}{c}{\textbf{Exceed limit if possible}} & \\\cmidrule(lr){2-3}
\textbf{Age} & Always & Not always & \textbf{Total}\\\midrule
Under 30 & 100 & 100 & 200\\ Over 30 & 40 & 160 & 200\\\midrule Total & 140 & 260 & 400\\\bottomrule
\end{tabular}
\end{center}
\end{context}
\begin{question}{50}
{Among the people with age under 30, what are the odds that they always exceed the speed limit?\par}
\optsfour{1 to 2}{2 to 1}{1 to 1}{50\%}
\answer{C}{1 to 1}
\begin{explain}
Odds $=$ always : not always $=100:100=\mathbf{1\ to\ 1}$.
\tip{Odds compare \emph{for} with \emph{against}; risk (probability) compares \emph{for} with the \emph{total}.}
\end{explain}
\end{question}

\begin{question}{51}
{What is the \% difference in the positive attitude of over 30 people towards not exceeding speed limit?\par}
\optsfour{60\%}{25\%}{40\%}{30\%}
\answer{D}{30\%}
\begin{explain}
The positive attitude is ``not always'' exceeding. Over 30: $160/200=80\%$; under 30: $100/200=50\%$. Difference $=\mathbf{30\%}$.
\end{explain}
\end{question}

\begin{question}{52}
{Overall, what is the "risk" of always exceeding the speed limit?\par}
\optsfour{0.33}{0.35}{0.65}{0.5}
\answer{B}{0.35}
\begin{explain}
Overall risk $=\dfrac{140}{400}=\mathbf{0.35}$.
\end{explain}
\end{question}

\begin{question}{53}
{What is the relative risk of always exceeding the speed limit for people under 30 compared to people over 30?\par}
\optsfour{2.5}{0.4}{0.5}{30\%}
\answer{A}{2.5}
\begin{explain}
Risk under 30 $=100/200=0.5$; risk over 30 $=40/200=0.2$.
\[\text{Relative risk}=\frac{0.5}{0.2}=\mathbf{2.5}\]
\end{explain}
\end{question}

\begin{question}{54}
{Among the people with age over 30, what is the "risk" of always exceeding the speed limit?\par}
\optsfour{0.20}{0.40}{0.33}{0.50}
\answer{A}{0.20}
\begin{explain}
Risk over 30 $=\dfrac{40}{200}=\mathbf{0.20}$.
\end{explain}
\end{question}

\begin{question}{55}
{Open learning system means:\par}
\optsone{student can take up any number of programmes at a time}{student can join a programme at any time}{student can join any programme at any time}{student can appear in the examination at any time}
\answer{C}{student can join any programme at any time}
\begin{explain}
In an \textbf{open learning} system entry barriers are removed. Per the key, the student \textbf{can join any programme at any time}, with flexibility in admission, pace, place and choice of courses.
\end{explain}
\end{question}

\begin{question}{56}
{`Impact factor' in research is linked with:\par}
\optstwo{Published books}{Research journals}{Research papers}{Research thesis}
\answer{B}{Research journals}
\begin{explain}
The \textbf{impact factor} (Eugene Garfield; Clarivate's Journal Citation Reports) is a measure for \textbf{journals}: the average number of citations received in a year by articles the journal published in the previous two years.
\end{explain}
\end{question}

\begin{question}{57}
{The birthday of late Dr. Abul Kalam Azad is celebrated as \blank{} in India.\par}
\optstwo{Teachers' Day}{National Higher Education Day}{National Education Day}{National Education Policy Day}
\answer{C}{National Education Day}
\begin{explain}
Maulana Abul Kalam Azad's birthday, \textbf{11 November}, is celebrated as \textbf{National Education Day} (since 2008).
\end{explain}
\end{question}

\begin{question}{58}
{In the credit based grading system (CGS) the term `credit' refers to the number of hours assigned for:\par}
\optstwo{a course per day}{a course per month}{a course per week}{the entire course}
\answer{C}{a course per week}
\begin{explain}
In a credit system, \textbf{one credit $=$ one hour of teaching (lecture or tutorial) per week} over a semester, or two hours of practical per week.
\end{explain}
\end{question}

\begin{question}{59}
{Private universities are being established in Gujarat. These come into existence on the basis of:\par}
\optstwo{Public demand}{State Legislative Act}{Government order}{UGC Directive}
\answer{B}{State Legislative Act}
\begin{explain}
A private university in India is established by an \textbf{Act of the State Legislature} (in Gujarat, under the Gujarat Private Universities Act, 2009) and must conform to the UGC (Private Universities) Regulations, 2003.
\end{explain}
\end{question}

\begin{question}{60}
{NAAC accreditation of a university reveals:\par}
\optstwo{the level of quality in the university}{the merits of the university}{the weaknesses of university}{the advice given to the university}
\answer{A}{the level of quality in the university}
\begin{explain}
NAAC accreditation certifies the \textbf{level of quality} of an institution, expressed as a grade (A++ to C) based on the CGPA.
\end{explain}
\end{question}
